\providecommand{\main}{..}		% Override relative path to the main file (already set in main file)
\documentclass[../InterneDSLs.tex]{subfiles}
\begin{document}

\chapter{Domänenspezifische Bedarfe}
Domänenspezifische Sprachen sind formale Sprachen, die für eine bestimmte Domäne (Fachbereich) entworfen wurden. Durch diese Modellierung soll eine hohe Spezifität im Bezug auf die Probleme der Domäne gewährleistet werden; und daneben soll nichts anderes ausgedrückt werden können.

Dadurch soll die Sprache für Domänenexperten leicht zu verwenden sein, im Gegensatz zu einer General Purpose Language (siehe~\ref{FIG:DslAdvantages}). Daraus ergibt sich auch, dass der Code besser lesbar ist und weniger Redundanz enthält.

Bei der Entwicklung einer eigenen DSL läuft man aber auch Gefahr, dass diese nur in einer kleinen Nische verwendet wird. Dadurch könnte es auftreten, dass mit vielen DSLs gearbeitet wird, die jeweils nur ihre eigene Domäne abdecken, und man unter Umständen viele Sprachen lernen muss (\ref{FIG:DslDisadvantages}).

\begin{figure}[ht]
\centering
  \begin{subfigure}[c]{0.49\textwidth}
    \begin{itemize}
	  \item Domänenexperten entwickeln mit
	  \item Lesbarkeit, weniger Redundanz
	  \item Änderung des Ausführungskontextes
    \end{itemize}
    \caption{Vorteile von DSLs}
    \label{FIG:DslAdvantages}
  \end{subfigure}
  \begin{subfigure}[c]{0.49\textwidth}
    \begin{itemize}
	  \item Viele fragmentierte DSLs
	  \item Implementierungsaufwand
	  \item Eigene DSL in kleiner Nische
    \end{itemize}
    \caption{Nachteile von DSLs}
    \label{FIG:DslDisadvantages}
  \end{subfigure}
  \caption{Vor- und Nachteile von DSLs}
  \label{FIG:DslAdvantagesDisadvantages}
\end{figure}


\section{Domäne}
Die Domäne kann aus der Informatik selbst kommen oder eine fremde Domäne sein. Beispiele aus der Informatik sind Build-Systeme (Make, Maven, ...) und SQL.

\section{Terminologie}

\subsubsection{Hostsprache}
\subsection{Konkrete Syntax}
\subsection{Abstrakte Syntax}
\subsection{Statische Semantik}
\subsection{Dynamische Semantik}

\section{Probleme und Anwendungsfälle}


\section{Interne und Externe Umsetzung}
Domänenspezifische Sprachen lassen sich auf zwei Arten kategorisieren. Zuerst kann man sie in externe und interne DSLs einteilen, und zum Anderen kann man sie in ihrer Notation unterscheiden.

\subsection{Externe DSLs}
Externe DSLs sind eigenständige, neu entwickelte Sprachen. Die \gls{Konkrete Syntax} und Semantik können selbst festgelegt werden und gelten deshalb als mächtiger. Allerdings muss mehr Aufwand für die Erstellung der Tools (Parser, Interpreter/Compiler, \acs{IDE}-Unterstützung) aufgewandt werden. Beispiele für externe DSLs sind LaTeX, Awk und SQL.

\subsection{Interne DSLs}
Interne DSLs werden innerhalb einer Hostsprache definiert und können deshalb dessen Tools mitbenutzen, wodurch der Implementierungsaufwand sinkt. Je nach verwendeter Hostsprache ähneln interne DSLs eher einer Bibliothek oder wirken wie eine eigene Sprache. Da der Aufwand bei internen DSLs geringer ist, eignen sie sich besser als Prototyp in agilen Umgebungen.~\cite{butting2018deriving}

Interne DSLs können zum Beispiel als Fluent-\acs{API} umgesetzt werden.


\chapter{Grammatiken}
Programmiersprachen liegen formale Sprachen zugrunde, die sich in der Chomsky-Hie"-rar"-chie (Abbildung~\ref{fig:chomkskyhierarchie}) wiederfinden. Diese formalen Sprachen legen durch Regeln fest, welche Zeichenketten gültig sind und damit zur Sprache gehören. Jede Klasse von Grammatiken beinhaltet auch die Grammatiken tiefer in der Chomsky-Hierarchie. (Überführung ineinander?)

\begin{figure}
\centering
\includegraphics[width=0.5\textwidth]{\main/10_Pictures/Grammatiken}
\caption{Chomsky-Hierarchie}
\label{fig:chomkskyhierarchie}
\end{figure}

\section{Typ-3-Grammatiken}\label{sec:refulaeregrammatik}
Typ-3-Grammatiken entsprechen regulären Grammatiken, die kontextfreie Grammatiken (siehe Abschnitt~\ref{sec:kontextfreiegrammatik}) mit speziellem Einschränkungen sind. Es gibt links- und rechtsreguläre Grammatiken, jedoch ist mit \gqq{regulär} meistens \gqq{rechtsregulär} gemeint. Bei linksregulären Grammatiken sind auf der rechten Seite der Bildungsregeln Nicht-Terminale nur links von den Terminalen erlaubt; bei rechtsregulären Grammatiken dürfen Nicht-Terminale entsprechend nur rechts von den Terminalen stehen. Ein Beispiel für reguläre Grammatiken sind reguläre Ausdrücke.

\section{Typ-2-Grammatiken}\label{sec:kontextfreiegrammatik}
Kontextfreie Grammatiken erlauben auf der rechten Seite der Bildungsregeln eine beliebige Folge von Terminalen und Nicht-Terminalen; auf der linken Seite darf nur ein Nicht-Terminal stehen, das also keinen Kontext hat. Es sind aber mehrere Bildungsregeln für dasselbe Nicht-Terminal erlaubt. Die meisten Programmiersprachen verwenden kontextfreie Grammatiken.\footcite{Fowler.2010}

\section{Typ-1-Grammatiken}\label{sec:kontextsensitivegrammatik}
Kontextsensitive Grammatiken erlauben auf der linken Seite der Bildungsregeln links und rechts von dem Nicht-Terminal weitere Terminal- und Nicht-Terminal-Symbole. Dadurch können diese Nicht-Terminale nur im Kontext zu anderen Symbolen auftreten. Auch sollte die rechte Seite der Bildungsregel mindestens gleich lang sein wie die linke. Analog wie bei Regulären Grammatiken (siehe Abschnitt~\ref{sec:refulaeregrammatik}) gibt es links-kontextsensitive und rechts-kontextsensitive Grammatiken, bei denen links respektive rechts vom Nicht-Terminal-Symbol weitere Symbole stehen.

\section{Typ-0-Grammatiken}\label{sec:beliebigegrammatik}
Beliebige formale Grammatik haben keine Einschränkungen, außer dass bei keiner Bildungsregel auf der linke Seite die leere Zeichenkette stehen darf. Der einzige Unterschied zu kontextsensitiven Grammatiken ist der, dass das Nicht-Terminal-Symbol neben einem Kontext auf der rechten Seite auch der leere String werden darf.


\section{Notation}
DSLs können auch in der Art der Notation unterschieden werden, es gibt grafische und textuelle DSLs. Ein Beispiel für eine grafische DSL ist UML, textuelle DSLs sind zum Beispiel HTML oder LaTeX. Fowler stellt auch Workbenches vor.\footcite[][S. 22ff]{Fowler.2010}

\subsection{EBNF}

\subsection{Grafisch}

\end{document}
